\lyseis{7.7}{Δεκαδική μορφή ρητών}
\sol{1} α) $0{,}375$ \quad β) $0{,}8\overline{3}$ \quad γ) $0{,}\overline{36}$ \quad δ) $0{,}\overline{7}$ \quad ε) $0{,}45$ \quad στ) $0{,}1\overline{3}$ \quad ζ) $0{,}08\overline{3}$.
\sol{2} Ακριβείς: $\frac{7}{20}$ ($20=2^2\cdot5$), $\frac{11}{40}$ ($40=2^3\cdot5$), $\frac{3}{16}$ ($16=2^4$), $\frac{13}{25}$ ($25=5^2$). Περιοδικοί: $\frac{5}{12}$ και $\frac{9}{14}$.
\sol{3} $0{,}\overline{36}=0{,}36363636$, $2{,}1\overline{45}=2{,}14545454$, $0{,}\overline{142857}=0{,}14285714$.
\sol{4} α) $\frac{7}{9}$ \quad β) $\frac{15}{99}=\frac{5}{33}$ \quad γ) $\frac{4}{3}$ \quad δ) $\frac{81}{99}=\frac{9}{11}$ \quad ε) $\frac{7}{30}$.
\sol{5} α) $0{,}\overline{3}>0{,}33$ \quad β) $0{,}\overline{6}=\frac{2}{3}$ \quad γ) $0{,}7\overline{2}=0{,}7222\ldots>0{,}72$.
\sol{6} α) $0{,}67$ \quad β) $0{,}83$ \quad γ) $0{,}14$ \quad δ) $0{,}44$.
\sol{7} α) Σωστό. β) Λάθος: $0{,}\overline{3}$. γ) Σωστό. δ) Σωστό. ε) Σωστό. στ) Λάθος: $0{,}1\overline{6}$.
\sol{8} Η περίοδος είναι $142857$. $2:7=0{,}\overline{285714}$ και $3:7=0{,}\overline{428571}$: ίδια ψηφία με το $1:7$, ξεκινούν από άλλο σημείο.
\sol{9} Το $10:3=3{,}\overline{3}$. Στην πράξη το καθένα παίρνει $3{,}33$ € και μένει $1$ λεπτό.
\sol{10} $\frac{100}{12}=\frac{25}{3}=8{,}\overline{3}$ m/s $\approx8{,}33$ m/s.
\sol{11} $\frac{1}{3}=33{,}\overline{3}\%\approx33{,}3\%$ και $\frac{5}{6}=83{,}\overline{3}\%\approx83{,}3\%$.
\sol{12} $\frac{23}{3}=7{,}\overline{6}\approx7{,}67$.
\sol{13} Τα υπόλοιπα της διαίρεσης με το $7$ είναι $1,\dots,6$. Αφού υπάρχουν μόνο έξι διαφορετικά, μετά από το πολύ έξι βήματα επαναλαμβάνεται κάποιο και μαζί του και τα ψηφία.
\sol{14} $x=0{,}\overline{9}$, $10x=9{,}\overline{9}$, $10x-x=9$, $9x=9$, $x=1$.
\sol{15} $\frac{1}{11}=0{,}\overline{09}$ και $\frac{1}{13}=0{,}\overline{076923}$.
\sol{16} Για παράδειγμα $0{,}335$, αφού $0{,}333\ldots<0{,}335<0{,}34$.
\sol{17} Το $0{,}333$ είναι προσέγγιση. Αν τα $3$ της περιόδου συνεχίζονται για πάντα, $3\cdot0{,}\overline{3}=0{,}\overline{9}=1$, όπως $3\cdot\frac{1}{3}=1$.
\sol{18} Ακριβή αναπτύγματα έχουν τα $\frac{1}{2},\frac{1}{4},\frac{1}{5},\frac{1}{8},\frac{1}{10}$ (και $\frac{1}{16}$ κ.λπ.). Οι παρονομαστές τους έχουν μόνο παράγοντες $2$ και $5$.
